% !TEX root = ../../ctfp-print.tex

\lettrine[lhang=0.17]{幺}{半群是范畴}理论和编程中都非常重要的概念。范畴对应强类型语言，
幺半群对应无类型语言。这是因为在幺半群中你可以组合任意两个箭头（就像在无类型语言中
可以组合任意两个函数——当然，运行程序时可能会遇到运行时错误）。

我们已经看到，幺半群可以被描述为具有单个对象的范畴，其中所有逻辑都编码在态射
组合规则中。这种范畴模型与更传统的集合论定义完全等价：在后者中我们将集合的两个元素
"相乘"得到第三个元素。这个"乘法"过程可以进一步分解为先形成元素对，然后将该元素对
与现有元素进行标识——即它们的"积"。

当我们放弃乘法的第二部分——将元素对与现有元素标识——会发生什么？我们可以从一个
任意集合开始，生成所有可能的元素对并称其为新元素。接着将这些新元素与所有可能的元素
配对，以此类推。这会产生链式反应——我们会无限添加新元素。最终得到的无限集合将
\emph{几乎}成为幺半群。但幺半群还需要单位元和结合律。没问题，我们可以添加特殊的
单位元并通过标识某些元素对来满足单位元和结合律要求。

让我们通过简单示例来看这个过程。假设初始集合有两个元素 $\{a, b\}$，我们将其称为
自由幺半群的生成元。首先添加特殊元素 $e$ 作为单位元。接着添加所有元素对并称其为
"积"。$a$ 和 $b$ 的积就是元素对 $(a, b)$，$b$ 和 $a$ 的积是 $(b, a)$，
$a$ 与 $a$ 的积是 $(a, a)$，$b$ 与 $b$ 的积是 $(b, b)$。我们也可以形成与 $e$
的配对如 $(a, e)$、$(e, b)$ 等，但会将它们分别标识为 $a$、$b$ 等。因此本轮只需添加
$(a, a)$、$(a, b)$、$(b, a)$、$(b, b)$，最终得到集合
$\{e, a, b, (a, a), (a, b), (b, a), (b, b)\}$。

\begin{figure}[H]
  \centering
  \includegraphics[width=0.8\textwidth]{images/bunnies.jpg}
\end{figure}

\noindent
下一轮将继续添加类似 $(a, (a, b))$、$((a, b), a)$ 等元素。此时需要确保持有结合律，
因此要将 $(a, (b, a))$ 与 $((a, b), a)$ 进行标识。换句话说，我们不再需要内部括号。

可以预见这个过程的最终结果：我们将创建所有可能的 $a$ 和 $b$ 序列。事实上，如果将 $e$
表示为空序列，就能发现我们的"乘法"本质上就是列表拼接。

这种持续生成所有可能组合并通过最小数量的标识来维持定律的构造方式被称为自由构造。
我们刚刚完成的就是从生成元集合 $\{a, b\}$ 构造出的 \newterm{自由幺半群}。

\section{Haskell中的自由幺半群}

Haskell中的二元集合对应类型 \code{Bool}，由该集合生成的自由幺半群对应类型
\code{{[}Bool{]}（布尔列表）。这里刻意忽略了无限列表带来的问题。）

Haskell通过类型类定义幺半群：

\src{snippet01}
这表明每个 \code{Monoid} 必须包含中性元素 \code{mempty} 和二元运算 \code{mappend}。
单位元和结合律无法在Haskell中直接表达，程序员每次实例化幺半群时都需要手动验证。

任意类型的列表构成幺半群可通过如下实例定义：

\src{snippet02}
该定义声明空列表 \code{{[}{]}} 为单位元，列表拼接运算符 \code{(++)} 为二元运算。

正如所见，类型为 \code{a} 的列表对应以 \code{a} 为生成元的自由幺半群。自然数集合配备
乘法运算不是自由幺半群，因为我们标识了大量乘积。例如对比：

\src{snippet03}
虽然这很直观，但问题是：能否在不允许观察对象内部结构的范畴论中执行这种自由构造？
我们将使用得力工具——泛构造来解决这个问题。

第二个有趣的问题是：任何幺半群都能通过对某个自由幺半群增加超出定律要求的最小标识数量
而获得吗？我将向你展示这可以直接从泛构造推出。

\section{自由幺半群的泛构造}

如果你还记得我们之前关于泛构造的经历，你可能会注意到它与其说是关于构造某个对象，不如说是选择最适合特定模式的对象。因此如果我们想用泛构造来“构造”一个自由幺半群，我们必须考虑从一大类幺半群中进行挑选。我们需要一个完整的幺半群范畴来进行选择。但幺半群真的能构成一个范畴吗？

首先让我们将幺半群视为带有额外结构的集合，这些结构由单位元和乘法定义。我们选择那些保持幺半群结构的函数作为态射。这种保持结构的函数被称为\newterm{同态映射}。一个幺半群同态必须将两个元素的乘积映射为各自映射后的乘积：

\src{snippet04}
并且必须将单位元映射到单位元。

例如，考虑从整数列表到整数的同态映射。如果我们把\code{{[}2{]}}映射为2，把\code{{[}3{]}}映射为3，就必须把\code{{[}2, 3{]}}映射为6，因为连接操作

\src{snippet05}
会转化为乘法

\src{snippet06}
现在让我们暂时忽略单个幺半群的内部结构，仅将其视为具有对应态射的对象。这样就得到了幺半群范畴$\cat{Mon}$。

不过，在忘记内部结构之前，让我们先注意一个重要性质。$\cat{Mon}$中的每个对象都可以平凡地映射到一个集合——正是它的元素集合。这个集合被称为\newterm{底层}集合。事实上，不仅对象可以映射到集合，$\cat{Mon}$的态射（同态映射）也可以映射为函数。虽然这看似平凡，但很快就会变得有用。这种从$\cat{Mon}$到$\Set$的对象和态射的映射实际上是一个函子。由于这个函子「遗忘」了幺半群结构——一旦进入普通集合范畴，我们就不再区分单位元或关心乘法运算——它被称为\newterm{忘却函子}。忘却函子在范畴论中经常出现。

现在我们有了$\cat{Mon}$的两种不同视角。我们可以像对待其他范畴一样处理它，这时看不到幺半群的内部结构。对$\cat{Mon}$中的特定对象，我们只能说它通过态射与其他对象（包括自身）相连。态射的「乘法表」——即组合规则——源自另一种视角：将幺半群视为集合。通过进入范畴论，我们并没有完全失去这种视角——仍然可以通过忘却函子访问它。

要应用泛构造，我们需要定义一个特殊性质，使我们能在幺半群范畴中搜索并选出自由幺半群的最佳候选。但自由幺半群是由其生成元定义的。不同的生成元选择会产生不同的自由幺半群（\code{Bool}列表与\code{Int}列表不同）。我们的构造必须始于一组生成元。于是我们又回到了集合范畴！

这就是忘却函子发挥作用的地方。我们可以用它来透视我们的幺半群。我们可以在这些「斑点」的X光图像中识别生成元。具体如下：

我们从一组生成元$x$开始——这是$\Set$中的一个集合。

我们要匹配的模式包含一个幺半群$m$（$\cat{Mon}$的对象）和一个$\Set$中的函数$p$：

\src{snippet07}
其中$U$是我们从$\cat{Mon}$到$\Set$的忘却函子。这是一个奇特的异构模式——一半在$\cat{Mon}$，一半在$\Set$。

这个想法是：函数$p$将在$m$的X光图像中识别生成元集合。即使函数可能无法精确标识集合中的点（可能产生坍缩），这也没关系。泛构造会通过选择最佳代表来解决这个问题。

\begin{figure}[H]
  \centering
  \includegraphics[width=0.4\textwidth]{images/monoid-pattern.jpg}
\end{figure}

\noindent
我们还需要定义候选对象间的排序关系。假设有另一个候选：一个幺半群$n$和一个在其X光图像中识别生成元的函数：

\src{snippet08}
当存在一个幺半群态射（保持结构的同态映射）

\src{snippet09}
使得其在$U$下的像（记住，$U$是函子，所以它将态射映射为函数）能够通过$p$分解时：

\src{snippet10}
我们就说$m$优于$n$。如果将$p$视为在$m$中选择生成元；将$q$视为在$n$中选择「相同」生成元；那么$h$就可以看作是在两个幺半群间映射这些生成元。根据定义，$h$保持幺半群结构意味着一个幺半群中两个生成元的乘积会被映射为另一个幺半群中对应生成元的乘积，依此类推。

\begin{figure}[H]
  \centering
  \includegraphics[width=0.4\textwidth]{images/monoid-ranking.jpg}
\end{figure}

\noindent
这种排序可用于寻找最佳候选——自由幺半群。定义如下：

\begin{quote}
  当且仅当对于任何其他幺半群$n$（连同函数$q$），存在唯一的态射$h$从$m$出发满足上述分解性质时，我们称$m$（连同函数$p$）是生成元为$x$的\textbf{自由幺半群}。
\end{quote}
顺便提一下，这回答了我们的第二个问题。$U h$就是那个有能力将$U m$的多个元素坍缩为$U n$单个元素的函数。这种坍缩对应于识别自由幺半群中的某些元素。因此，任何以$x$为生成元的幺半群都可以通过对基于$x$的自由幺半群进行某些元素识别得到。自由幺半群正是只进行了必要最少识别的那个幺半群。

当我们讨论邻接时会再次回到自由幺半群的概念。

\section{挑战}

\begin{enumerate}
  \tightlist
  \item
        你可能会认为（如同我最初的想法一样），幺半群同态保持单位元的要求是多余的。毕竟我们已知对所有 $a$ 都有：

        \begin{snip}{text}
h a * h e = h (a * e) = h a
\end{snip}
        因此 $h e$ 表现得像一个右单位元（类似地也表现为左单位元）。问题在于，对所有 $a$ 而言，$h a$ 可能仅覆盖目标幺半群的一个子幺半群。在 $h$ 的像之外可能存在真正的单位元。证明：若两个幺半群之间保持乘法的同构映射必然自动保持单位元。
  \item
        考虑从整数列表（带连接(concatenation)运算）到整数（带乘法(multiplication)运算）的幺半群同态映射。空列表 \code{{[}{]}} 的像是什么？假设所有单元素列表都被映射为其包含的整数，即 \code{{[}3{]}} 映射到 3 等。那么 \code{{[}1, 2, 3, 4{]}} 的像是什么？有多少个不同的列表映射到整数 12？这两个幺半群之间是否存在其他同态？
  \item
        单元素集合生成的自由幺半群(free monoid)是什么？你能看出它与什么结构同构吗？
\end{enumerate}